
\section{A Little from Our District Meeting}

“Folkebladet,” November 29, 1916

Pastor Joh. Skonnord began with the sermon. His subject was the Christian hope. Among other things, he said that it was a great riches to possess the hope that does not put to shame, because God’s love has been poured out in our hearts.

The subject for discussion, Luke 10:23–37, was introduced by Pastor Gustav Sandanger of Vattle Læke. It was easy to understand that the speaker had devoted much labor to the introductory address he gave, and it was especially well suited to lead us into the thoughts of the text. To “introduce” a word of God, or any other subject for Christian discussion, is a very difficult matter, which few are able to accomplish. I myself have felt the difficulty very keenly when I have had the task of giving such an introduction. Pastor Sandanger made a very good attempt at introducing the word of God before us. The sum of his thoughts concerning this passage of Scripture seemed to me capable of being formulated thus: Neither belonging to God’s people, nor having religious experiences, nor possessing orthodox views of the way of salvation avails unto salvation, but a life whose task is to show mercy.

When Jesus spoke these words, he had before his mind’s eye a people of God who were zealous for good works.

During the discussion of this word of God, several emphasized that the condition for coming into possession of a Samaritan disposition was to be permitted to “see Jesus”: “Blessed are the eyes that see what you see,” Jesus says to his disciples.

During the discussion, in which many of the pastors present took part, the importance of being permitted to “see Jesus” was emphasized with great force: to see him in the way the disciples saw him, as the Savior of the world, their Savior.

If I were to point to something that seemed to me to have been neglected at the district meeting, I would mention that little was said about the man who had fallen among robbers. In other words, it seemed to me that too little was said about mission, home mission, and mission to the heathen. There was, however, one of the pastors who spoke briefly, in passing, about home or inner mission. He emphasized how important it is for our congregations to support this work and to make it possible to labor with the Word of God among our people who live in the newly settled districts and who there are exposed to the danger of forgetting their soul and their God. It is my firm conviction that this district meeting will bear some fruit for several of those who attended it.

Since the secretary is publishing a complete report of the meeting, I shall pass over the details. I believe it would be a great blessing to the life of the congregations if every district in the Free Church, at one of its district meetings, could take up our common undertakings: Augsburg, the missions, and the work of mercy.

E. B.